\begin{lema}[Rank of the line equalities]\label{lem:rango-llegada}
A realizable type of \(v\) distinct sites specifies their four partitions into row, column, diagonal, and
antidiagonal fibers. Let \(c_s\) be its number of components and \(D\) the dimension of the space of
positions. If component \(g\) contains \(v_g\) sites and has cyclomatic number \(\beta_g\), its rank is
\(v_g-1+L_g\), with \(L_g\ge\min(\beta_g,2)\). In particular, \(D=v+c_s-\sum_gL_g\le p-1\) for the three
exponents in Proposition~\ref{prop:soportes-llegada}.
\end{lema}
\begin{proof}
Replace each fiber by a tree on its sites. Its edges impose equalities \(n_f\cdot(z_i-z_l)=0\), where
\(n_f\) is one of the four normals. This preserves the equalities, the components, and \(E-V+c_s\). A
spanning tree of a component gives \(v_g-1\) independent equations by leaf elimination. A dependence on a
subgraph would assign nonzero coefficients to its edges so that at each vertex the signed sum of the
weighted normals vanishes. Its support has no leaves. At a vertex of degree two both edges lie on the
same line. Suppress each path of degree-two vertices. A component consisting solely of such vertices, or
a loop after suppression, would form a cycle in the tree of a line. Two parallel edges would give two
straight paths between the same two distinct sites, which determine a unique line. They too would close a
cycle in that tree. The reduced graph is simple and has minimum degree three. In each component,
\(V\ge4\) and \(E-V+1\ge V/2+1\ge3\). Suppressing the paths preserves the number of cycles. Thus a
spanning tree with one or two chords contains the support of no dependence. Adding \(\min(\beta_g,2)\)
chords gives the stated rank. Subtracting this rank from \(2v\) and summing over the components gives the
last inequality.
\end{proof}

\begin{lema}[Weights and number of types]\label{lem:tipos-llegada}
Set \(L_s=v+c_s-D\) and \(d_i=\max_{f\ni i}|f\cap U|\). For fixed \(v,c_s,L_s\) the sum over realizable
labeled types satisfies
\[
 \frac1{v!}\sum_{\rm tipos}\prod_i d_i\le \frac{C^v(4v^2)^{L_s}}{c_s!}.
\]
Suppose that the only diagonal conflicts in the support are two pairs on distinct lines. With four sites,
the bounds \((D,p)\) for zero through four axial coincidences are \((6,7),(5,6),(4,5),(3,4),(3,4)\) when
the class is realizable. With three sites the pairs share one site, and the bounds are \((4,5),(3,4)\)
according as their endpoints have zero or one axial coincidence.
\end{lema}
\begin{proof}
For each site choose the first fiber on which it attains its maximum \(d_i\) in a fixed ordering of the
families. Join the \(a_f>0\) sites assigned to each fiber \(f\) to a node of its color. Add equality
edges between sites until each component is connected, without forming cycles. The stars admit an
equivalent tree of equalities between sites. This forest has rank \(v-c_s\). Completing the basis
requires \(L_s\) edges, each with at most \(4v^2\) choices. The basis recovers the four partitions, since
an equality in its span holds in the realization and therefore belongs to the type. Root each component
at one site and join the roots under a common node. There are at most \(2v+1\) nodes, a fixed number of
colors, and \(v\) site labels. The plane tree encoding has at most \(C^v v!\) possibilities by the
Catalan bound. Each type admits at least \(c_s!\prod_f(a_f-1)!\) distinct orderings of the children. The
members of each star are identified by their labels. With \(D_f=|f\cap U|\),
\[
 \prod_f\frac{D_f^{a_f}}{(a_f-1)!}\le 2^v\exp\!\left(\sum_fD_f\right)
 \le 2^ve^{4v}.
\]
This gives the weighted sum. When counting unordered supports, each support gives \(v!\) labelings, even
when several have the same type.

The classification of the two conflicts includes all axial coincidences. By hypothesis, the two pairs
account for all diagonal edges. With no axial coincidences, the four sites form two trees. The first
coincidence joins them, the second creates a cycle, and from the third onward the general bound gives
\(p=4\). The rank lemma reduces the dimension to \(6,5,4\) in the first three cases. With three
coincidences the two rows and one column of the second pair are fixed given those of the first, or
conversely. At most three parameters remain, even if there is a fourth coincidence. In the three-site
wedge the initial graph is a tree and an axial coincidence between its endpoints creates a cycle. If
there is a third
diagonal conflict, the cycles are counted again in the full graph. The rank lemma and
Proposition~\ref{prop:soportes-llegada} then give \(D\le p-1\) with the exponent of that graph. Pairs in
the same square are eliminated by subtracting its error.
\end{proof}
